% Part: second-order-logic
% Chapter: metatheory
% Section: introduction

\documentclass[../../../include/open-logic-section]{subfiles}

\begin{document}

\olfileid{sol}{met}{int}

\olsection{سريزه}

د لومړۍ درجې منطق څو ښې ځانګړنې لري۔ پوهېږو چې د پرېکړې وړ نۀ دے،
خو لږ تر لږه د بديهي اصولو په بڼه وړاندې کېدے شي۔ يعنې د لومړۍ درجې
منطق دپاره داسې ثبوتي نظامونه شته چې سم او بشپړ دي: د !!a{derivability}
اړيکه~$\Proves$ داسې ټاکي چې د !!{sentence}s د هر سټ~$\Gamma$ او هرې
!!{sentence}~$!A$
دپاره $\Gamma \Entails !A$ هله او يوازې هله چې $\Gamma \Proves !A$۔
له دې په ځانګړې توګه دا راځي چې د لومړۍ درجې منطق معتبرې جملې
!!{computably
  enumerable} دي۔ يوه محاسبه کېدونکې تابع~$f\colon \Nat \to
\Sent[L]$ شته چې د~$f$ ارزښتونه د~$\Lang{L}$ ټولې او يوازې معتبرې
!!{sentence}s دي۔ ځکه !!{derivation}s شمېرلے شو او هغه اشتقاقونه چې يوه
واحده !!{sentence} !!{derive} کوي، هماغې !!{sentence} ته لېږو۔ د دويمې درجې منطق
د لومړۍ درجې له منطق څخه د بيان زيات ځواک لري؛ نو د هغه د معتبرو
جملو راټولول عموماً لا پېچلي دي۔ په حقيقت کښې به وښيو چې د دويمې
درجې منطق يوازې ناپرېکړه‌وړ نۀ دے، بلکې معتبرې جملې ئې حتا
!!{computably
  enumerable} هم نۀ دي۔ له دې امله د دويمې درجې منطق
دپاره سم او بشپړ ثبوتي نظام نۀ شي کېدے (که څه هم سم خو نابشپړ
ثبوتي نظامونه شته او په څېړنه کښې مهم دي)۔

د لومړۍ درجې منطق دوه نورې ځانګړنې هم لري: کمپکت دے (که د !!{sentence}s
د سټ~$\Gamma$ هر متناهي فرعي سټ د صدق وړ وي، نو $\Gamma$ پخپله هم
د صدق وړ دے)؛ او د لوېنهايم--سکولم قضيه پکې صدق کوي (که $\Gamma$
نامتناهي مدل ولري، نو !!a{denumerable} مدل هم لري)۔ دا دواړه پايلې
د دويمې درجې منطق دپاره ناسمې دي۔ لامل بيا هم دا دے چې د دويمې درجې
منطق د !!{domain}s د اندازې په باب داسې واقعيتونه بيانولے شي چې د
لومړۍ درجې منطق ئې نۀ شي بيانولے۔

\end{document}
